% ch56_en_algebras.tex — Volume IV Chapter 20
% Retrofitted with full Vol I-III pedagogy: cycle stages, etymology, dialogs.

\chapter{$\mathbb{E}_n$-Algebras and the Little Disks Operads}

\begin{overviewbox}
The \term{little $n$-disks operad} $\mathbb{E}_n$ is the most important
family of operads in higher algebra. Its $n$-th component
$\mathbb{E}_n(k)$ is the space of configurations of $k$ disjoint
$n$-disks inside the unit $n$-disk. Algebras over $\mathbb{E}_n$ in a
symmetric monoidal $\infty$-category — \term{$\mathbb{E}_n$-algebras} —
interpolate between fully associative ($\mathbb{E}_1$) and fully
commutative ($\mathbb{E}_\infty$) structures, indexed by codimension.

This chapter introduces:
\begin{itemize}
  \item Configuration spaces $\mathrm{Conf}_k(\mathbb{R}^n)$ and the
        action of $\Sigma_k$.
  \item The little disks operad $\mathbb{E}_n$ in $\mathrm{Top}$ (or
        $\mathcal{S}$ after fibrant replacement).
  \item $\mathbb{E}_n$-algebras in symmetric monoidal $\infty$-categories.
  \item The endpoints: $\mathbb{E}_1 = \mathrm{Assoc}$, $\mathbb{E}_\infty = \mathrm{Comm}$.
  \item \term{Dunn additivity}: $\mathbb{E}_m \otimes \mathbb{E}_n \simeq
        \mathbb{E}_{m+n}$ — the operadic tensor product of little disks
        interpolates dimensions additively.
\end{itemize}

The applications cited (not developed in detail): factorization homology,
commutative ring spectra, structured ring spectra, the cohomology of
loop spaces (May's recognition theorem). The reader who has completed
this chapter can open Lurie's \emph{Higher Algebra} Chapter~5 directly.
\end{overviewbox}


% ═══════════════════════════════════════════════════════════════════════════
\section{Configuration Spaces}
\label{sec:configurations}
% ═══════════════════════════════════════════════════════════════════════════

\subsection{Formalizing}

\begin{definition}[Configuration space]
\label{def:configuration}
For $n \geq 0$ and $k \geq 0$, the \term{configuration space} is
\begin{equation*}
  \mathrm{Conf}_k(\mathbb{R}^n) \;=\; \bigl\{(x_1, \ldots, x_k) \in (\mathbb{R}^n)^k : x_i \neq x_j \text{ for } i \neq j\bigr\}.
\end{equation*}
The symmetric group $\Sigma_k$ acts on $\mathrm{Conf}_k(\mathbb{R}^n)$ by
permuting coordinates.
\end{definition}

\begin{example_thm}[Low-dimensional examples]
\label{ex:conf-low}
$\mathrm{Conf}_2(\mathbb{R}) = \{(x, y) \in \mathbb{R}^2 : x \neq y\}$ has
two connected components (separated by $x = y$). Modulo $\Sigma_2$, one
component.

$\mathrm{Conf}_2(\mathbb{R}^2) \simeq \mathbb{R}^4 \setminus \mathrm{diag}$
has homotopy type of $S^1$ (the angle between two points). Quotient by
$\Sigma_2$: real projective line, $\mathbb{RP}^1 \simeq S^1$.

$\mathrm{Conf}_k(\mathbb{R}^\infty)$ is contractible for every $k$ —
the infinite-dimensional configuration space ``has room for everything.''
\end{example_thm}


% ═══════════════════════════════════════════════════════════════════════════
\section{The Little $n$-Disks Operad $\mathbb{E}_n$}
\label{sec:little-disks}
% ═══════════════════════════════════════════════════════════════════════════

\subsection{Formalizing}

\begin{nameorigin}
\term{Little disks operad} was introduced by \emph{J.\ Michael Boardman}
and \emph{Rainer Vogt} (1973) in their book \emph{Homotopy Invariant
Algebraic Structures on Topological Spaces}. The original
operad-of-little-cubes models the structure of $n$-fold loop spaces:
the space of $k$ disjoint $n$-cubes embedded in a larger $n$-cube
parameterizes ``$k$-fold operations'' that can be performed on an
$n$-fold loop space. Boardman--Vogt's recognition theorem characterizes
$n$-fold loop spaces as algebras over the $\mathbb{E}_n$ operad.

The notation \term{$\mathbb{E}_n$} is due to \emph{Peter May}: ``E'' for
\emph{everything} (an $\mathbb{E}_\infty$-algebra has \emph{everything}
commutative, including all higher homotopies), with the subscript $n$
recording the dimension at which commutativity is enforced. The naming
convention has become universal in higher algebra.
\end{nameorigin}

\begin{definition}[Little disks $\mathbb{E}_n$]
\label{def:little-disks}
For $1 \leq n \leq \infty$, the \term{little $n$-disks operad}
$\mathbb{E}_n$ has:
\begin{itemize}
  \item $\mathbb{E}_n(k) = $ the space of \term{rectilinear embeddings}
        of $k$ disjoint $n$-disks into the unit $n$-disk. (A rectilinear
        embedding scales and translates.)
  \item $\Sigma_k$-action by permuting the disks.
  \item Composition: nesting embeddings.
\end{itemize}
We have $\mathbb{E}_n(k) \simeq \mathrm{Conf}_k(\mathbb{R}^n)$ (the
configuration of disk centers, weakly equivalently).
\end{definition}

\begin{howtosay}
$\mathbb{E}_n$ \sayit{``E-sub-n'' or ``the little $n$-disks operad''};
also written $\mathbb{E}_n$ or $E_n$ depending on the typographic
convention. $\mathbb{E}_\infty$ is the colimit as $n \to \infty$.
\end{howtosay}

\begin{theorem}[$\mathbb{E}_n$ as an $\infty$-operad]
\label{thm:en-infty-operad}
$\mathbb{E}_n$ has a canonical structure of an $\infty$-operad in
$\mathcal{S}$ (the $\infty$-category of spaces), with its $k$-th
component being the (weak homotopy type of the) configuration space.
\end{theorem}


% ═══════════════════════════════════════════════════════════════════════════
\section{$\mathbb{E}_n$-Algebras and the Endpoints}
\label{sec:en-algebras}
% ═══════════════════════════════════════════════════════════════════════════

\subsection{Formalizing}

\begin{definition}[$\mathbb{E}_n$-algebra]
\label{def:en-algebra}
For a symmetric monoidal $\infty$-category $\mathcal{C}^\otimes$, an
\term{$\mathbb{E}_n$-algebra in $\mathcal{C}$} is an algebra over the
$\infty$-operad $\mathbb{E}_n$:
\begin{equation*}
  \mathrm{Alg}_{\mathbb{E}_n}(\mathcal{C}) \;=\; \mathrm{Fun}_{\mathrm{Fin}_*}^{\mathrm{inert}}(\mathbb{E}_n^\otimes, \mathcal{C}^\otimes).
\end{equation*}
\end{definition}

\begin{theorem}[Endpoints of the $\mathbb{E}_n$-tower]
\label{thm:en-endpoints}
\begin{itemize}
  \item $\mathbb{E}_1$-algebras in $\mathcal{C}^\otimes$ are
        \term{associative monoids in $\mathcal{C}^\otimes$} (the
        coherent version of associative algebras).
  \item $\mathbb{E}_\infty$-algebras in $\mathcal{C}^\otimes$ are
        \term{commutative algebras in $\mathcal{C}^\otimes$} (coherent
        commutative algebras).
\end{itemize}
For each $1 \leq n \leq \infty$, $\mathbb{E}_n$ interpolates:
$\mathbb{E}_n$-algebras have $n$-disk-parametrized commutativity, weaker
than full commutativity for $n < \infty$ and weaker than associativity
for $n = 0$ (which trivializes).
\end{theorem}

\begin{example_thm}[Ring spectra]
\label{ex:ring-spectra}
$\mathbb{E}_n$-algebras in $\mathrm{Sp}$ are \term{$\mathbb{E}_n$-ring
spectra}. Most prominent cases:
\begin{itemize}
  \item $\mathbb{E}_1$-ring spectrum $=$ associative ring spectrum
        (Lurie's ``$\mathbb{A}_\infty$-ring spectrum'').
  \item $\mathbb{E}_\infty$-ring spectrum $=$ commutative ring spectrum.
        Underlies most modern stable homotopy theory: $\mathbb{S}$ (the
        sphere spectrum), $H\mathbb{Z}$, $K$-theory $KU$, complex
        cobordism $MU$, all are $\mathbb{E}_\infty$.
\end{itemize}
\end{example_thm}

\begin{example_thm}[$\mathbb{E}_2$ in topology: braided monoids]
\label{ex:e2-braided}
$\mathbb{E}_2$-algebras in $\mathbf{Set}$ (or $\mathcal{S}$): braided
monoidal structures. The $\mathbb{E}_2$-operad's configuration spaces
have non-trivial $\pi_1$ (braids), encoding the data of \emph{half-twists}
between commuting elements.
\end{example_thm}


% ═══════════════════════════════════════════════════════════════════════════
\section{Dunn Additivity}
\label{sec:dunn-additivity}
% ═══════════════════════════════════════════════════════════════════════════

\subsection{Formalizing}

\begin{nameorigin}
\term{Dunn additivity} is named for \emph{Gerald Dunn}, who proved
the theorem in his 1988 paper ``Tensor product of operads and iterated
loop spaces.'' The result is striking: the operadic tensor product
$\mathbb{E}_m \otimes \mathbb{E}_n \simeq \mathbb{E}_{m+n}$ shows that
codimensions \emph{add} under the tensor product of little-disks
operads. The Eckmann-Hilton argument is the $m = n = 1$ special case:
two compatible monoid structures on a set are equivalent to a single
commutative monoid (i.e., $\mathbb{E}_1 \otimes \mathbb{E}_1 \simeq
\mathbb{E}_2$, weakened in the $1$-categorical case to commutativity).
\end{nameorigin}

\begin{theorem}[Dunn additivity]
\label{thm:dunn}
For any $m, n \geq 1$, there is an equivalence of $\infty$-operads
\begin{equation*}
  \mathbb{E}_m \otimes \mathbb{E}_n \;\simeq\; \mathbb{E}_{m+n},
\end{equation*}
where $\otimes$ is the Boardman--Vogt tensor product of $\infty$-operads.

Equivalently: $\mathbb{E}_{m+n}$-algebras are the same as
$\mathbb{E}_m$-algebras whose underlying multiplication is itself
$\mathbb{E}_n$-coherent — codimensions add.
\end{theorem}

\begin{example_thm}[$\mathbb{E}_2 = \mathbb{E}_1 \otimes \mathbb{E}_1$]
\label{ex:e2-e1-tensor}
A doubly-monoidal structure (an $\mathbb{E}_1$-algebra in
$\mathbb{E}_1$-algebras) is exactly an $\mathbb{E}_2$-algebra. This is the
$\infty$-categorical version of the Eckmann--Hilton argument (Chapter~19):
two compatible monoid structures form a braided monoid.
\end{example_thm}

\begin{example_thm}[Limit case]
\label{ex:e-infty-limit}
$\mathbb{E}_\infty = \mathrm{colim}_n \mathbb{E}_n$, achieved by stabilizing
the configuration spaces. Dunn additivity gives $\mathbb{E}_\infty \otimes
\mathbb{E}_n \simeq \mathbb{E}_\infty$ — adding any finite $n$ to infinity
remains infinity.
\end{example_thm}


% ═══════════════════════════════════════════════════════════════════════════
\section{Exercises}
\label{sec:en-exercises}
% ═══════════════════════════════════════════════════════════════════════════

\begin{exercisesbox}
\begin{qa}
\qaitem{Define $\mathrm{Conf}_k(\mathbb{R}^n)$.}{$k$-tuples of distinct points in $\mathbb{R}^n$. $\Sigma_k$ acts by permutation.}
\qaitem{What is $\mathrm{Conf}_2(\mathbb{R}^2)$ homotopy-equivalent to?}{$S^1$ (the angle between the two points, modulo translations and scalings).}
\qaitem{What is $\mathrm{Conf}_k(\mathbb{R}^\infty)$?}{Contractible — every configuration can be ``spread out'' in the infinite-dimensional space.}
\qaitem{Define the little $n$-disks operad $\mathbb{E}_n$.}{$\mathbb{E}_n(k) = $ rectilinear embeddings of $k$ disjoint $n$-disks into the unit $n$-disk. Equivalently (up to weak equivalence), $\mathrm{Conf}_k(\mathbb{R}^n)$ with $\Sigma_k$-action.}
\qaitem{What's $\mathbb{E}_\infty$?}{$\mathrm{colim}_n \mathbb{E}_n$: the limit as $n \to \infty$. All configurations become equivalent because $\mathrm{Conf}_k(\mathbb{R}^\infty)$ is contractible.}
\qaitem{Define an $\mathbb{E}_n$-algebra.}{Algebra over the $\mathbb{E}_n$-operad in a symmetric monoidal $\infty$-category. Concretely: an object with $n$-disk-parametrized multiplication.}
\qaitem{What's an $\mathbb{E}_1$-algebra?}{Associative monoid (up to coherent homotopy). The most familiar coherent algebra structure.}
\qaitem{What's an $\mathbb{E}_\infty$-algebra?}{Commutative monoid (coherently). All higher coherences for commutativity are explicit data.}
\qaitem{Why do $\mathbb{E}_n$-structures interpolate between associativity and commutativity?}{Because configuration spaces $\mathrm{Conf}_k(\mathbb{R}^n)$ have increasingly trivial fundamental groups as $n$ grows: $\pi_1 = $ braid group for $n=2$, becomes simpler for $n \geq 3$.}
\qaitem{Give the principal example in $\mathrm{Sp}$.}{$\mathbb{E}_\infty$-ring spectra: $\mathbb{S}$ (sphere), $H\mathbb{Z}$, $KU$, $MU$. All foundational ring spectra of stable homotopy theory.}
\qaitem{What's an $\mathbb{E}_2$-algebra in $\mathbf{Set}$?}{A braided monoidal set: commutative-up-to-braid-twist; rare in concrete sets, common in higher categorical settings.}
\qaitem{State Dunn additivity.}{$\mathbb{E}_m \otimes \mathbb{E}_n \simeq \mathbb{E}_{m+n}$. Tensor product of little-disks operads adds dimensions.}
\qaitem{Why is Dunn additivity natural?}{Because configurations in $\mathbb{R}^{m+n}$ ``decompose'' as configurations in $\mathbb{R}^m$ times configurations in $\mathbb{R}^n$ (in the limit of nested compatibility).}
\qaitem{What does ``$\mathbb{E}_m$-algebra in $\mathbb{E}_n$-algebras'' produce?}{By Dunn: an $\mathbb{E}_{m+n}$-algebra. Eckmann--Hilton in disguise: doubly-monoidal $=$ braided $=$ $\mathbb{E}_2$.}
\qaitem{Give a non-trivial example of Dunn at work.}{$\mathbb{E}_3$ as $\mathbb{E}_1 \otimes \mathbb{E}_2$: three commuting monoid structures, with the inner two compatible via a braid. Higher coherence between three operations.}
\qaitem{Are all $\mathbb{E}_\infty$-rings (in $\mathrm{Sp}$) commutative?}{Coherently commutative, yes. Strict commutativity is rare and typically forced by being concentrated in degree 0; coherent commutativity allows higher data without being strict.}
\qaitem{What's factorization homology?}{An $\mathbb{E}_n$-algebra invariant for $n$-manifolds: $\int_M A$ for $M$ an $n$-manifold and $A$ an $\mathbb{E}_n$-algebra. Generalizes Hochschild homology ($n=1$, $S^1$).}
\qaitem{What's May's recognition theorem?}{An $\mathbb{E}_n$-algebra in $\mathcal{S}_*$ is equivalent to an $n$-fold loop space $\Omega^n X$ for some $X$. The classical recognition theorem, generalized to all $n$ and to all symmetric monoidal $\infty$-cats.}
\qaitem{What's the $\mathbb{E}_n$-center?}{For an $\mathbb{E}_n$-algebra $A$, the center $Z_{\mathbb{E}_n}(A)$ is the largest $\mathbb{E}_{n+1}$-algebra mapping to $A$. Generalizes the classical center of an associative algebra.}
\qaitem{Where does this chapter take us?}{Foundations of higher algebra à la Lurie's HA Chapter 5. Reader is prepared to encounter $\mathbb{E}_n$-algebras in actual contexts: ring spectra, derived geometry, factorization homology.}
\qaitem{What's the next chapter?}{Ch~57: $A_\infty$ and $L_\infty$ algebras. Strong-homotopy associative ($A_\infty$, equivalent to $\mathbb{E}_1$) and Lie ($L_\infty$, Koszul dual of $A_\infty$) algebras, with explicit operadic presentations.}
\qaitem{Are $\mathbb{E}_n$-operads the only operads in higher algebra?}{No — $\mathrm{Lie}_\infty$, $A_\infty$, $C_\infty$, and many bespoke operads play key roles. The $\mathbb{E}_n$-family is foundational because of Dunn additivity.}
\qaitem{What's the symmetric counterpart of $\mathbb{E}_n$?}{$\mathbb{E}_\infty$ for fully symmetric (commutative). The $\Sigma_k$-action on $\mathbb{E}_n(k)$ becomes increasingly homotopy-trivial as $n$ grows.}
\end{qa}
\end{exercisesbox}


% ═══════════════════════════════════════════════════════════════════════════
\section*{Chapter Summary}
\addcontentsline{toc}{section}{Chapter Summary}
% ═══════════════════════════════════════════════════════════════════════════

\begin{summarybox}
\textbf{Configuration spaces.}\quad $\mathrm{Conf}_k(\mathbb{R}^n)$:
$k$-tuples of distinct points in $\mathbb{R}^n$. Foundational object;
$\mathbb{E}_n(k) \simeq \mathrm{Conf}_k(\mathbb{R}^n)$ as $\Sigma_k$-spaces.

\textbf{Little disks $\mathbb{E}_n$.}\quad $\infty$-operad whose
$k$-th component is the configuration of $k$ disjoint $n$-disks in
the unit $n$-disk.

\textbf{$\mathbb{E}_n$-algebra.}\quad An algebra over $\mathbb{E}_n$ in
a symmetric monoidal $\infty$-category. $\mathbb{E}_1 = $ associative;
$\mathbb{E}_\infty = $ commutative; intermediate $n$ interpolate.

\textbf{Foundational examples.}\quad In $\mathrm{Sp}$: $\mathbb{S}, H\mathbb{Z},
KU, MU$ are all $\mathbb{E}_\infty$-ring spectra. $\mathbb{E}_2$ in
$\mathcal{S}$: braided monoidal spaces (e.g., $2$-fold loop spaces).

\textbf{Dunn additivity.}\quad $\mathbb{E}_m \otimes \mathbb{E}_n \simeq
\mathbb{E}_{m+n}$. Codimensions add under operadic tensor product.
Eckmann--Hilton as a special case.

\textbf{Bridge to applications.}\quad Factorization homology, ring
spectra, May's recognition theorem, $\mathbb{E}_n$-center,
$\mathbb{E}_n$-Hochschild homology.
\end{summarybox}


% ═══════════════════════════════════════════════════════════════════════════
\section*{Formal Restatement}
\addcontentsline{toc}{section}{Formal Restatement}
% ═══════════════════════════════════════════════════════════════════════════

\begin{formalrestatement}
\textbf{Configuration space.}\quad
$\mathrm{Conf}_k(\mathbb{R}^n) = \{(x_1, \ldots, x_k) \in (\mathbb{R}^n)^k : x_i \neq x_j\}$.

\medskip
\textbf{$\mathbb{E}_n$-operad.}\quad
$\mathbb{E}_n(k) \simeq \mathrm{Conf}_k(\mathbb{R}^n)$ as $\Sigma_k$-space;
composition by nesting.

\medskip
\textbf{$\mathbb{E}_n$-algebra.}\quad
$\mathrm{Alg}_{\mathbb{E}_n}(\mathcal{C})$: $\infty$-functor
$\mathbb{E}_n^\otimes \to \mathcal{C}^\otimes$ preserving inert coCartesian.

\medskip
\textbf{Endpoints.}\quad
$\mathbb{E}_1 = \mathrm{Assoc}$; $\mathbb{E}_\infty = \mathrm{Comm}$.

\medskip
\textbf{Dunn additivity.}\quad
$\mathbb{E}_m \otimes \mathbb{E}_n \simeq \mathbb{E}_{m+n}$ for $m, n \geq 0$.

\medskip
\textbf{Examples.}\quad
$\mathbb{E}_\infty$-ring spectra: $\mathbb{S}, H\mathbb{Z}, KU, MU$;
$\mathbb{E}_2$ in $\mathcal{S}$: braided $\infty$-monoids;
$\mathbb{E}_n$-algebras in $\mathrm{Sp}$ for general $n$.
\end{formalrestatement}
